% Source: https://artofproblemsolving.com/community/c3h2529353_chinese_remainder_theorem
% Stable identity: aops-topic-2529353-q01
\textbf{(a)} \emph{Lemma.} If each $m_i$ divides an integer $N$, then $M\mid N$. Indeed, every prime power $p^e$ exactly dividing some $m_i$ divides $N$, and since the $m_i$ are pairwise relatively prime, different $m_i$ involve different primes. So the prime factorisation of $M$ divides that of $N$.

\emph{Uniqueness.} If $x$ and $y$ are both solutions, then each $m_i$ divides $x-y$, so $M\mid x-y$ by the lemma. Conversely, anything congruent to a solution modulo $M$ is again a solution.

\emph{Existence.} Send each $r\in\{0,1,\dots,M-1\}$ to the $k$-tuple of its remainders $\left(r\bmod m_1,\dots,r\bmod m_k\right)$. If $r\ne s$ gave the same tuple, then $M\mid r-s$ by the lemma, which is impossible since $0<|r-s|<M$. So this map is injective. It goes from a set of $M$ elements to the set of all tuples of remainders, which also has $m_1\cdots m_k=M$ elements, so it is a bijection. In particular, some $r$ has remainders $\left(a_1\bmod m_1,\dots,a_k\bmod m_k\right)$, and $x=r$ is a solution. $\blacksquare$

\textbf{(b)} The conditions $x\equiv2\pmod3$ and $x\equiv2\pmod7$ together say $x\equiv2\pmod{21}$ (by (a), or directly since $3$ and $7$ both divide $x-2$). Among $2,23,44,65,86$ (one full set modulo $105$), the one that is $\equiv3\pmod5$ is $23$. So the solutions are $x\equiv23\pmod{105}$, and the smallest positive one is $\boxed{23}$.
